\documentclass[11pt,a4paper]{article}
\usepackage[T1]{fontenc}
\usepackage[utf8]{inputenc}
\usepackage{times}
\usepackage[a4paper,margin=2.5cm]{geometry}
\usepackage{amsmath,amssymb}
\usepackage{graphicx}
\usepackage{booktabs}
\usepackage{array}
\usepackage{rotating}
\usepackage[font=small,labelfont=rm]{caption}
\usepackage[switch,modulo]{lineno}
\usepackage{enumitem}
\usepackage[obeyspaces,spaces]{url}
\makeatletter\g@addto@macro{\UrlBreaks}{\do\_\do\=\do\-\do\.}\makeatother
\usepackage[hidelinks]{hyperref}
\usepackage{microtype}
\usepackage{seqsplit}
\modulolinenumbers[5]
\newcommand{\Msun}{M_\odot}
\DeclareUrlCommand\file{\urlstyle{tt}}
\newcommand{\ee}[2]{#1\times10^{#2}}
\hypersetup{pdftitle={The Time Attraction Law (TAL), Release V21},pdfauthor={David Szabolcsi}}

\title{\textbf{The Time Attraction Law (TAL), Release V21:\\
Conditional Neutron-Star Equilibrium Sequences\\
with a Locally Centred Top-Hat Density Average}}
\author{David Szabolcsi\\[2pt]
\small Independent experimental researcher, Unterterzen, Switzerland\\[2pt]
\small ORCID: 0009-0005-6300-074X}
\date{}

\begin{document}
\maketitle
\linenumbers

\begin{abstract}
\noindent The Time Attraction Law (TAL) postulates that the local rate of proper time depends on the
top-hat averaged density $\bar\rho_L$ through $f_L=\exp(-\gamma_L\bar\rho_L)$ with $\gamma_L=GL^2/c^2$.
This paper reports a conditional sensitivity study of compact stars in which the top-hat average is
centred at every stellar radius and recomputed self-consistently on the equilibrium solution. Two
labelled ans\"atze insert $f_L$ either into the mass equation (``mass-source'') or into the pressure
and lapse gradients (``gradient'') of the Tolman--Oppenheimer--Volkoff system, for the SLy, APR4 and H4
piecewise-polytropic fits and $L=1$, 367, 1000, 10\,000 and 100\,000~m: 33 equilibrium sequences in total.
A finite-source bound shows that only the neutron-star regime can reach a percent-level TAL factor
(3.52--5.79\,\% for compactness 0.15--0.25); white dwarfs and the Sun cannot. At fixed gravitational
mass the largest radius shift is 0.164~km at $1.4\,\Msun$, against a three-EOS GR spread of 2.632~km,
and 0.623~km at $2.0\,\Msun$, against 1.557~km; in this comparison the radius effect is EOS-degenerate.
All 1.4 and $2.0\,\Msun$ comparison stars lie in the causal range of the fits. APR4 becomes superluminal
($c_s^2=dP/d\epsilon>c^2$) above $\rho_0=\ee{1.4458}{15}$~g\,cm$^{-3}$ already in GR; maxima beyond this
density are reported only as labelled, potentially acausal benchmarks, with causal endpoints given
separately. The turning-point criterion $dM/d\rho_c>0$ is applied as a heuristic only. Under an additional,
separately assumed additive clock coupling, Sirius~B conditionally excludes a finite stellar-scale band
$\ee{1.117}{7}\lesssim L\lesssim\ee{3.477}{7}$~m, but imposes no monotonic upper limit on $L$ and does not
constrain the 1~m--100~km values. NICER results are shown as an illustrative overlay of GR-derived
measurements, not as a TAL likelihood. No detection is claimed, no unconditional bound on $L$ is derived,
and the coupling-consistency task 50aa2db0 remains open.
\end{abstract}

\noindent\textbf{Keywords:} time attraction law; density-dependent proper time; top-hat averaged density;
neutron stars; equation of state; causality; NICER; Sirius~B

\section{Introduction}
In general relativity the rate of an ideal clock is set by the metric \cite{will}. The Time Attraction Law
(TAL) \cite{tal} is a phenomenological proposal in which the local proper-time rate is instead an explicit
function of the density averaged over a sphere of radius $L$. The previous release (Zenodo V20 \cite{v20})
studied a laboratory coil configuration, and the release before it (V19 \cite{v19}) the TAL factor as a
seventh model for SPARC rotation curves. The V20 manuscript (\file{V18/manuscript/tal_v20_fop.pdf} in \cite{gh}) reports a calculated laboratory
signal of $\approx-\ee{1.20}{-48}$ and a SPARC TAL factor numerically identical to the pure-baryon baseline
at the published scales.
A natural question is therefore where, if anywhere, a top-hat TAL factor can become large.

The purpose of this paper is narrow and deliberately conservative. It (i) states a finite-source bound
showing that among ordinary stars, white dwarfs and neutron stars only the neutron-star regime can reach a
percent-level factor; (ii) documents exactly which modified equilibrium equations were solved, with a true,
locally centred top-hat; (iii) reports maximum masses and fixed-mass radii for 33 conditional sequences and
compares the radius shifts with the spread of three standard equations of state (EOSs); (iv) flags the
superluminal high-density part of the EOS fits and reports results in their causal range; (v) applies the
turning-point criterion as a heuristic; and (vi) gives a conditional Sirius~B redshift consistency check.
All numbers are taken from the archived V21 package (Zenodo V21 \cite{v21}; GitHub V19, directory \file{V21/} \cite{gh});
file names are given so that each value can be traced.

The central result is a set of conditional sensitivity calculations. The two insertions of $f_L$ are
ans\"atze, not solutions of a completed TAL field theory, and the backlog consistency task 50aa2db0 may
change which insertion, if any, is admissible. Nothing in this paper is a detection, an unconditional bound
on $L$, or a closed demonstration of physical consistency.

\section{Theoretical framework}
\subsection{Definitions}
The definitions are those of the TAL formalism document \file{Time_Attraction.pdf} \cite{tal}, as quoted in the
V20 manuscript (\file{V18/manuscript/tal_v20_fop.pdf} in \cite{gh}; Def.~1--3, Eqs.~(1), (2), (7) of \cite{tal}). The top-hat averaged density at a point $\mathbf x$ is
\begin{equation}
\bar\rho_L(\mathbf x)=\frac{3}{4\pi L^3}\int_{|\mathbf x'-\mathbf x|\le L}\rho(\mathbf x')\,d^3x',
\end{equation}
the coupling constant is
\begin{equation}
\gamma_L=\frac{G L^2}{c^2},\qquad [\gamma_L]=\mathrm{m^3\,kg^{-1}},
\end{equation}
and the time factor is
\begin{equation}
f_L(\bar\rho_L)=\exp(-\gamma_L\bar\rho_L),\qquad d\tau_L=f_L\,dt,\qquad f_L\in(0,1].
\end{equation}
The source gives an upper bound $L<367$~m from a generic clock-stability argument (Eq.~(3), p.~3 of \cite{tal});
it is the author's own estimate, not a constraint from a dedicated experiment, and it is not imposed here.
The value $L=367$~m is included in the scan as one of the requested scales. At $L=10$~km the constants of
the package give $\gamma=\ee{7.4262}{-20}$~m$^3$\,kg$^{-1}$ (\file{note_HU_EN.md} \S1).

\subsection{Finite-source bound}
Write $x_L=\gamma_L\bar\rho_L$, $f_L=e^{-x_L}$. For positive density and a Euclidean averaging ball,
$x_L(\mathbf x)=3G M_{\mathrm{ball}(\mathbf x,L)}/(4\pi c^2 L)$. For a spherically symmetric, radially
decreasing density, a ball of given radius contains most mass when centred at the stellar centre, so
\begin{equation}
\sup_{\mathbf x,L}x_L(\mathbf x)\le\frac{3}{4\pi}\max_L\frac{G M_0(<L)}{Lc^2},\qquad
0\le1-f_L\le1-e^{-x_{\max}}\le x_{\max},
\end{equation}
where $M_0(<L)=4\pi\int_0^L\rho(s)s^2\,ds$ is the unweighted coordinate-density integral; in the mass-source
model it is not the modified gravitational mass $m(L)$ (\file{note_HU_EN.md} \S1). Inserting compactness scales
gives the orders of magnitude of Table~\ref{tab:bound}. These are scale estimates, not rigorous ceilings:
central concentration can increase $\max GM_0(<L)/(Lc^2)$ relative to the surface value. For a finite source
$x_L\propto L^2$ at small $L$ in smooth interior matter, but $x_L\propto1/L$ once the whole source is
enclosed, so an ever larger $L$ does not create an arbitrarily large stellar correction. The statement
concerns finite stellar sources, not an extended homogeneous cosmological background.

\begin{table}[t]
\centering
\caption[]{Scale estimates of the maximal TAL exponent for finite stellar sources. Source: \file{note_HU_EN.md} \S1.}
\label{tab:bound}
\begin{tabular}{lccc}
\toprule
Object & Compactness scale $GM/(Rc^2)$ & $x\sim(3/4\pi)C$ & $1-e^{-x}$\\
\midrule
Neutron star & 0.15--0.25 & 0.0358--0.0597 & 3.52--5.79\,\%\\
Typical white dwarf & $10^{-4}$ & $\ee{2.39}{-5}$ & 0.00239\,\%\\
Sun & $\ee{2}{-6}$ & $\ee{4.77}{-7}$ & 0.0000477\,\%\\
\bottomrule
\end{tabular}
\end{table}

\subsection{Density convention and equilibrium prescriptions}
Let $\rho_0$ be the EOS rest-mass density, $\epsilon$ the total energy density and $\rho=\epsilon/c^2$. The
latter is used both in the mass equation and in the TAL averaging integral; this recovers GR at $f=1$.
Using $\rho_0$ instead would be a different model, not a notational substitution (\file{note_HU_EN.md} \S2).
With
\begin{equation}
A(r)=\frac{G\,[m(r)+4\pi r^3P(r)/c^2]}{c^2r^2\,[1-2Gm(r)/(rc^2)]}
\end{equation}
and the lapse notation $g_{tt}=-e^{2\nu}c^2$, the three prescriptions of Table~\ref{tab:presc} were
evaluated. The gradient alternative multiplies the gradient of $\nu$, not the lapse itself; no model applies
both modifications.

\begin{table}[t]
\centering
\caption[]{Equilibrium prescriptions. Labels as in the package files. Source: \file{note_HU_EN.md} \S2.}
\label{tab:presc}
\begin{tabular}{lll}
\toprule
Label & Mass equation & Pressure / lapse gradient\\
\midrule
\file{GR} & $m'=4\pi r^2\epsilon/c^2$ & $P'=-(\epsilon+P)A$, $\nu'=A$\\
\file{mass} (published mass-source insertion) & $m'=4\pi r^2(\epsilon/c^2)f_L$ & $P'=-(\epsilon+P)A$, $\nu'=A$\\
\file{gradient} (labelled alternative) & $m'=4\pi r^2\epsilon/c^2$ & $P'=-(\epsilon+P)f_LA$, $\nu'=f_LA$\\
\bottomrule
\end{tabular}
\end{table}

These are conditional ans\"atze. This paper does not derive a covariant action or supply missing effective
stresses. An extremum of a modified equilibrium sequence is not a demonstrated radial-stability boundary;
$M_{\max}$ below always means the maximum mass on the computed sequence.

\subsection{True top-hat, centred locally}
The density is zero outside the stellar surface $R$. At each stellar radius $r$ the averaging sphere is
centred at that point, not at the stellar centre. For $r>0$, full shells $0<s<L-r$, when present, contribute
$4\pi\int\rho s^2\,ds$; partial shells contribute
\begin{equation}
M_{\mathrm{cap}}(r,L)=\frac{\pi}{r}\int_{|r-L|}^{\min(r+L,R)}s\,\rho(s)\,[L^2-(s-r)^2]\,ds,
\end{equation}
with empty intervals omitted. At $r=0$, $M=4\pi\int_0^{\min(L,R)}\rho s^2\,ds$. The enclosed mass is divided
by the whole ball volume $4\pi L^3/3$, including any vacuum portion. Neither the central density nor
$m(r)/(4\pi r^3/3)$ is substituted for the local ball average. The convolution uses Euclidean coordinate
volume in areal coordinates, not a proper-volume or geodesic-ball definition; at neutron-star compactness
this choice is physically consequential and remains conditional (\file{note_HU_EN.md} \S2).

\section{Methodology}
\subsection{Equations of state}
The three EOSs are represented by the piecewise-polytropic fits of Read et al.\ \cite{read}, not by microscopic EOS
tables (Table~\ref{tab:eos}). Pressure and energy density are matched continuously, including the four-piece
SLy crust; the core boundaries are $10^{14.7}$ and $10^{15}$~g\,cm$^{-3}$. The low-density coefficients of
Read's Table~II are interpreted as $K/c^2$ in cgs mass-density units, and the later three crust coefficients
are fixed by continuity. The maximum relative error of the enthalpy inversion is $\ee{9.771}{-12}$ for each
EOS (\file{geometry_tests.json}).

\begin{table}[t]
\centering
\caption[]{Piecewise-polytropic parameters \cite{read}. Sources: \file{note_HU_EN.md} \S2; \file{geometry_tests.json}.}
\label{tab:eos}
\begin{tabular}{lcccc}
\toprule
EOS & $\log_{10}P_1$ (dyne\,cm$^{-2}$) & $\Gamma_1$ & $\Gamma_2$ & $\Gamma_3$\\
\midrule
SLy & 34.384 & 3.005 & 2.988 & 2.851\\
APR4 (AP4 in \cite{read}) & 34.269 & 2.830 & 3.445 & 3.348\\
H4 (stiff-radius comparison) & 34.669 & 2.909 & 2.246 & 2.144\\
\bottomrule
\end{tabular}
\end{table}

\subsection{Integration, self-consistency and scan}
Integration uses the logarithmic relativistic enthalpy $h=\int dP/(\epsilon+P)$, a central expansion at
$r=10^{-4}$~km and an outer endpoint $h=10^{-10}$. Each solution is iterated: solve the equilibrium
equations, rebuild the full density profile, recompute the top-hat, and repeat until the relative mass and
radius change and the maximum exponent change are all below $\ee{2}{-7}$. Production runs use ODE relative
tolerance $10^{-10}$, 801 radial profile points and 80-point Gauss shell quadrature
(\file{note_HU_EN.md} \S2; \file{stellar_tal.py}).

For every EOS, GR and both modifications were run for $L=1$, 367, 1000, 10\,000 and 100\,000~m: 33
sequences. The initial 53-point grid of central rest-mass density spans
$\log_{10}\rho_{0c}\,[\mathrm{g\,cm^{-3}}]=14.55$--15.65; a bounded local search refines each maximum.
Radii at fixed masses are found by root solving on the lower-central-density branch; a blank fixed-mass
entry means that the sequence cannot reach the target on that branch and is not extrapolated.

\subsection{Sound speed and causal endpoints}
\file{integrate()} evaluates $c_s^2/c^2=dP/d\epsilon$ ($=\Gamma P/(\epsilon+P)$ on each polytropic piece)
at every accepted ODE step, at every one of the 2401 output nodes and on both sides of every EOS piece
boundary (\file{note_HU_EN.md} \S7.2; \file{causality_stability_log.txt}). No causal completion of the EOS
was imposed. The ``causal endpoint'' of a sequence is the star whose centre sits exactly at the first
density with $c_s^2=c^2$; its mass is stored as \file{M_max_causal_Msun}. This is the endpoint of the
published fit's causal range on the followed branch, not the maximum mass of a causally completed EOS. The
areal radius of the $c_s^2=c^2$ core boundary is obtained from the threshold enthalpy $h_*$ (bracketed root,
\file{brentq}, xtol $=$ rtol $=10^{-15}$) evaluated on the dense output of the converged ODE solution;
the residual $|c_s^2/c^2-1|$ at $h_*$ is at most $\ee{2.2}{-16}$ (\file{note_HU_EN.md} \S7.4, H2).

\subsection{Turning-point heuristic}
The sign of $dM/d\rho_c$ was evaluated on the 53-point scan grid of all 33 sequences and, separately,
including the optimizer and root-finder points (\file{stellar_summary.csv}, columns
\file{dMdrho_sign_changes_scan_grid}, \file{dMdrho_sign_changes_all_points}). For the modified equations
this GR-motivated criterion is a heuristic, not a proven radial-stability result; the radial-mode
(Sturm--Liouville) analysis with TAL-modified perturbation equations has not been performed.

\subsection{NICER reference values}
The reference values of Table~\ref{tab:nicer} are read by \file{plots()} from
\file{observational_inputs.csv}; no NICER number is hard-coded. Uncertainties are marginal
16th--84th-percentile credible intervals. Each plotted reference is an uncorrelated, symmetrized Gaussian
approximation whose standard deviations are the arithmetic means of the published lower and upper errors;
ellipses satisfy $(\Delta R/\sigma_R)^2+(\Delta M/\sigma_M)^2=-2\ln(1-p)$ for joint $p=0.68$ and 0.95.
They are not the papers' posterior samples: mass--radius correlation, skewness and hot-spot multimodality
are absent. Before drawing, \file{check_observational_inputs()} verifies that each symmetrized error equals
the mean of the lower and upper errors (tolerance $10^{-9}$); all 10 source/quantity rows pass, with a largest
difference of $\ee{1.1}{-16}$ (\file{numerical_diagnostics.json}). RA/Dec (J2000) are metadata from the ATNF
Pulsar Catalogue v2.6.5 \cite{atnf}, cross-checked against SIMBAD, and are not used in the plot.

\begin{table}[t]
\centering
\caption[]{NICER reference values (marginal, approximately 68\,\%). The last two rows are kept for reference
and are not drawn (\file{used_in_plot=False}). Source: \file{observational_inputs.csv}; \file{note_HU_EN.md} \S4.}
\label{tab:nicer}
\begin{tabular}{llcc}
\toprule
Pulsar & Analysis & $M$ ($\Msun$) & $R$ (km)\\
\midrule
J0740+6620 & Salmi et al.\ 2024 \cite{salmi} & $2.073\pm0.069$ & $12.49^{+1.28}_{-0.88}$\\
J0437$-$4715 & Choudhury et al.\ 2024 \cite{choudhury} & $1.418\pm0.037$ & $11.36^{+0.95}_{-0.63}$\\
J0030+0451 & Kini et al.\ 2026, PDT-U (preferred) \cite{kini} & $1.43^{+0.20}_{-0.17}$ & $12.68^{+1.31}_{-1.04}$\\
\midrule
J0030+0451 & Kini et al.\ 2026, ST+PDT (alternative) \cite{kini} & $1.29^{+0.13}_{-0.12}$ & $11.63^{+0.95}_{-0.89}$\\
J0437$-$4715 & Miller et al.\ 2025, sensitivity \cite{miller} & $1.431\pm0.043$ & $13.892^{+1.25}_{-2.069}$\\
\bottomrule
\end{tabular}
\end{table}

All NICER posteriors were inferred with GR-based photon propagation, atmosphere models and rotational
geometry. Their comparison with nonrotating conditional TAL sequences is an illustrative overlay, not a
self-consistent alternative-gravity likelihood. For J0437$-$4715 the posterior mass error 0.037 is used,
not the radio-prior error 0.044. The J0030+0451 values are the six-year 2026 analysis \cite{kini}, not the older
2024 result. The Miller et al.\ reanalysis including modulated nonthermal emission \cite{miller} gives a bimodal radius
posterior with bounds 11.823 and 15.142~km; it is not combined with the Choudhury result.

\subsection{Sirius B: additive clock coupling}
Joyce et al.\ \cite{joyce} give the gravitational redshift $v_{\mathrm{gr}}=80.65\pm0.77$~km\,s$^{-1}$, the redshift
mass $1.017\pm0.025\,\Msun$ and $R=(0.00803\pm0.00011)\,R_\odot$; Bond et al.\ \cite{bond} give the independent
dynamical mass $M_B=1.018\pm0.011\,\Msun$. The dynamical mass and the independently estimated radius are used
to predict the GR redshift. The additional, clock-only assumption is
\begin{equation}
1+z=(1+z_{\mathrm{GR}})\exp(x_{\mathrm{em}}-x_{\mathrm{obs}}),\qquad
\delta v_B=c\,(1+z_{\mathrm{GR},B})\,[e^{x_B}-1].
\end{equation}
It retains a GR exterior and multiplies the local clock factor by TAL. It is not uniquely implied by either
equilibrium insertion of Table~\ref{tab:presc}; in particular, the gradient prescription does not itself
imply this multiplicative lapse law. The differential Sirius A/B strategy cancels the observer factor, and the
modelled Sirius~A contribution is subtracted. For $L\ge2R_B$ the ball around a surface emitter contains the
whole white dwarf, so that
\begin{equation}
x_B=\frac{3GM_B}{4\pi Lc^2},
\end{equation}
independently of the density profile. For $L<2R_B$, normalized Lane--Emden profiles with $n=1.5$ and $n=3$
at the observed mass and radius, with an $n=3$ Sirius~A profile, are used as bracketing illustrations, not
as a tabulated finite-temperature Sirius~B EOS (\file{note_HU_EN.md} \S5).

\subsection{Numerical checks}
Every EOS/coupling/$L$ case was re-solved at its peak central density and near $1.4\,\Msun$ with 1601
radial points, 160 quadrature points and ODE rtol $=10^{-11}$, giving 66 refinement comparisons
(\file{convergence.csv}). All 33 causal endpoints (14 of them below the sequence maximum) and the core
boundary of all 14 acausal maxima were recomputed with the same finer settings
(\file{causal_endpoint_convergence.csv}). A separate, polynomial-exact integral on the piecewise-linear profile
was compared with the 80-point production convolution in 15 geometry cases (all five $L$ at centre, half
radius and surface), and the code was tested against exact uniform-sphere limits, including the surface
overlap $\bar\rho_L/\rho=5/16$ for $L=R$ (\file{numerical_diagnostics.json}; \file{geometry_tests.json}).

\section{Results}
\subsection{GR baseline}
Table~\ref{tab:gr} lists the GR sequence maxima and fixed-mass radii. GR maxima and $R_{1.4}$ agree with the
fit predictions quoted by Read et al.\ \cite{read} to within 0.051\,\%; the small common offset is retained, not
renormalized away (\file{numerical_diagnostics.json}, \file{Read_GR_benchmarks}).

\begin{table}[t]
\centering
\caption[]{GR baseline. Sources: \file{stellar_summary.csv}; \file{causality_stability_log.txt}; \file{note_HU_EN.md} \S3.}
\label{tab:gr}
\begin{tabular}{lccc}
\toprule
EOS & Maximum sequence mass ($\Msun$) & $R(1.4\,\Msun)$ (km) & $R(2.0\,\Msun)$ (km)\\
\midrule
SLy & 2.048411 & 11.7061 & 10.6266\\
APR4 & 2.188001 & 11.3197 & 10.8667\\
H4 & 2.013851 & 13.9522 & 12.1834\\
\bottomrule
\end{tabular}
\end{table}

\subsection{Maximum masses and radii at $1.4\,\Msun$}
Table~\ref{tab:full} gives the full EOS $\times L$ comparison. All digits are numerical outputs, not
observational precision; differences at $L=1$~m are below the tested numerical error floor and their signs
have no physical meaning. All sequences pass the nominal $2\,\Msun$ mass-capability check before EOS causality
or modified radial stability is imposed; this is not evidence for TAL, since the GR baselines already pass it.
The mass-source insertion reduces the source at a fixed density profile, but the star readjusts, so its
maximum equilibrium mass need not decrease.

Only the APR4 GR sequence reaches the central J0740+6620 mass of $2.073\,\Msun$; the SLy 10~km modified
sequences also reach it, while H4 reaches 2.0599 (mass-source) and $2.0696\,\Msun$ (gradient). Given the
$0.069\,\Msun$ posterior uncertainty these statements are not EOS exclusions. The APR4 maxima lie above
the APR4 causal onset (Sect.~\ref{sec:causal}) and are therefore benchmarks of the published fit, not
physical mass limits.

\begin{table}[t]
\centering
\caption[]{Sequence maxima and radii at fixed gravitational mass $1.4\,\Msun$ (masses in $\Msun$, radii in km).
APR4 and SLy mass-source maxima at $L=1$, 10 and 100~km lie above the causal onset (Table~\ref{tab:causal}).
Sources: \file{stellar_summary.csv}; \file{note_HU_EN.md} \S3.}
\label{tab:full}
\begin{tabular}{lrcccc}
\toprule
EOS & $L$ (m) & $M_{\max}$ mass-source & $M_{\max}$ gradient & $R_{1.4}$ mass-source & $R_{1.4}$ gradient\\
\midrule
SLy & 1 & 2.048411 & 2.048411 & 11.7061 & 11.7061\\
SLy & 367 & 2.048591 & 2.048567 & 11.7065 & 11.7063\\
SLy & 1000 & 2.049739 & 2.049558 & 11.7091 & 11.7074\\
SLy & 10000 & 2.098219 & 2.108482 & 11.8703 & 11.7920\\
SLy & 100000 & 2.053290 & 2.058655 & 11.7356 & 11.7297\\
\midrule
APR4 & 1 & 2.188001 & 2.188001 & 11.3197 & 11.3197\\
APR4 & 367 & 2.188182 & 2.188160 & 11.3201 & 11.3199\\
APR4 & 1000 & 2.189339 & 2.189168 & 11.3227 & 11.3209\\
APR4 & 10000 & 2.241521 & 2.248766 & 11.4757 & 11.3937\\
APR4 & 100000 & 2.193438 & 2.198771 & 11.3471 & 11.3393\\
\midrule
H4 & 1 & 2.013851 & 2.013851 & 13.9522 & 13.9522\\
H4 & 367 & 2.013994 & 2.013975 & 13.9525 & 13.9523\\
H4 & 1000 & 2.014904 & 2.014767 & 13.9545 & 13.9532\\
H4 & 10000 & 2.059898 & 2.069632 & 14.0952 & 14.0267\\
H4 & 100000 & 2.018823 & 2.025597 & 13.9875 & 13.9821\\
\bottomrule
\end{tabular}
\end{table}

\subsection{Radius shifts versus EOS spread}
Table~\ref{tab:shift} compares the largest fixed-mass radius shift with the GR spread of the three EOSs. The
largest shifts occur at $L=10$~km: at $1.4\,\Msun$ for SLy (mass-source and gradient), at $2.0\,\Msun$ for H4
(both ans\"atze) (\file{radius_shifts.csv}). At $1.4\,\Msun$, $L=367$~m gives only 0.306--0.412~m of
mass-source radius increase and 0.139--0.178~m for the gradient ansatz; $L=1$~km gives about 2.26--3.04~m and
1.03--1.32~m; $L=100$~km is weaker again, approximately 27.4--35.3~m and 19.6--29.9~m. At $2.073\,\Msun$ only
one GR EOS reaches the mass, so \file{radius_shifts.csv} leaves the spread blank instead of reporting a
misleading zero; the APR4 10~km mass-source shift there, about 0.325~km, cannot be promoted into a multi-EOS
comparison.

\begin{table}[t]
\centering
\caption[]{Largest fixed-mass radius shift over the 30 modified sequences versus the three-EOS GR radius
spread. Sources: \file{radius_shifts.csv}; \file{note_HU_EN.md} \S3.}
\label{tab:shift}
\small
\begin{tabular}{cccc}
\toprule
Fixed $M$ ($\Msun$) & GR EOS spread (km) & Largest mass-source $\Delta R$ (km) & Largest gradient $\Delta R$ (km)\\
\midrule
1.4 & 2.6325 & 0.1642 & 0.0859\\
2.0 & 1.5569 & 0.6234 & 0.5475\\
\bottomrule
\end{tabular}
\end{table}

The chosen three-EOS range is a comparison envelope, not a probability interval or the full set of possible
EOSs. A shift smaller than this envelope is not a proof of exact degeneracy at every mass; it means that radius
alone does not identify TAL in this comparison. No EOS parameters or $L$ were fitted to NICER.

\subsection{Exponent profiles}
At $1.4\,\Msun$ and $L=10$~km the true centre exponents are 0.0321--0.0469, whereas extending the central
density uniformly across the averaging ball would give 0.0458--0.0738; the surface exponents are only
0.00616--0.01181 (\file{numerical_diagnostics.json}, entry \file{top_hat_profile_diagnostics}; also
\file{profiles_L10km.csv}). The factors vary through the star and are recomputed on the modified solution;
a central-density proxy would overstate them.

\subsection{Causality}\label{sec:causal}
Each EOS has a single acausal interval (Table~\ref{tab:onset}). For APR4, above
$\rho_0\approx\ee{1.446}{15}$~g\,cm$^{-3}$ ($\epsilon/c^2\approx\ee{1.801}{15}$~g\,cm$^{-3}$) the GR baseline
is already superluminal, and the modified branches inherit this. We report causal-range endpoints separately
while retaining the full equilibrium sequences as explicitly labelled, potentially acausal benchmarks. All 11
APR4 maximum-mass stars contain an acausal core whose $c_s^2=c^2$ boundary lies at $r=5.244$--5.411~km, with
$c_s^2/c^2$ up to 1.370--1.402 (Table~\ref{tab:causal}). The causal maximum mass is $2.0781\,\Msun$ in GR
(full sequence maximum $2.1880\,\Msun$) and 2.0781--$2.1414\,\Msun$ in the modified branches.

The SLy GR maximum is just causal ($c_s^2/c^2=0.99976$). The mass-source maxima at $L=1$, 10 and 100~km are
slightly acausal (up to 1.0112); the violation begins at 98.69--99.94\,\% of the central density and the
core boundaries lie at $r=0.237$, 1.067 and 0.230~km, respectively. Their causal maximum mass is at most
$\ee{1.5}{-4}\,\Msun$ below the sequence maximum. The gradient-ansatz SLy maxima and all H4 maxima are causal
(H4: $c_s^2/c^2\le0.522$). In all 33 sequences the 1.4 and $2.0\,\Msun$ comparison stars have central
densities below the onset, so every $R(1.4)$ and $R(2.0)$ in this paper, and the 0.164/2.632~km and
0.623/1.557~km comparison, lies in the causal range (\file{stellar_summary.csv}, \file{causal_at_1p4},
\file{causal_at_2p0}).

\begin{table}[t]
\centering
\caption[]{First superluminal density of each EOS fit. Sources: \file{causality_stability_log.txt};
\file{note_HU_EN.md} \S7.2.}
\label{tab:onset}
\small
\begin{tabular}{lccc>{\raggedright\arraybackslash}p{4.2cm}}
\toprule
EOS & onset $\rho_0$ (g\,cm$^{-3}$) & $\log_{10}\rho_0$ & onset $\epsilon/c^2$ (g\,cm$^{-3}$) & onset in \% of $\rho_{0c}$ at $M_{\max}$\\
\midrule
APR4 & $\ee{1.4458}{15}$ & 15.1601 & $\ee{1.8013}{15}$ & 74.49--76.28 (all 11 sequences)\\
SLy & $\ee{2.0071}{15}$ & 15.3026 & $\ee{2.8626}{15}$ & 98.69--99.94 (mass-source, $L=1$, 10, 100~km)\\
H4 & $\ee{1.0461}{16}$ & 16.0196 & $\ee{4.3012}{16}$ & never reached\\
\bottomrule
\end{tabular}
\end{table}

\begin{table}[t]
\centering
\small
\caption[]{Sequences whose maximum lies above the causal onset (14 of 33). $c_s^2$ is the maximum in the
$M_{\max}$ star; $r_{\mathrm{core}}$ is the areal radius of the $c_s^2=c^2$ boundary; $M_{\mathrm{caus}}$ and
$R_{\mathrm{caus}}$ refer to the causal endpoint. Source: \file{causality_stability_log.txt}.}
\label{tab:causal}
\footnotesize
\setlength{\tabcolsep}{4pt}
\begin{tabular}{llrcccccc}
\toprule
EOS & Mode & $L$ (m) & $M_{\max}$ ($\Msun$) & $c_s^2/c^2$ & onset (\%) & $r_{\mathrm{core}}$ (km) & $M_{\mathrm{caus}}$ ($\Msun$) & $R_{\mathrm{caus}}$ (km)\\
\midrule
APR4 & GR & -- & 2.188001 & 1.384439 & 75.463 & 5.246704 & 2.078144 & 10.6827\\
APR4 & mass & 1 & 2.188001 & 1.384303 & 75.470 & 5.246217 & 2.078144 & 10.6827\\
APR4 & mass & 367 & 2.188182 & 1.384654 & 75.451 & 5.247810 & 2.078266 & 10.6832\\
APR4 & mass & 1000 & 2.189339 & 1.385310 & 75.413 & 5.252319 & 2.079053 & 10.6870\\
APR4 & mass & 10000 & 2.241521 & 1.401687 & 74.491 & 5.411172 & 2.121668 & 10.9119\\
APR4 & mass & 100000 & 2.193438 & 1.385446 & 75.406 & 5.261463 & 2.082682 & 10.7217\\
APR4 & gradient & 1 & 2.188001 & 1.384070 & 75.484 & 5.245387 & 2.078144 & 10.6827\\
APR4 & gradient & 367 & 2.188160 & 1.384238 & 75.474 & 5.246427 & 2.078305 & 10.6827\\
APR4 & gradient & 1000 & 2.189168 & 1.384032 & 75.486 & 5.248500 & 2.079323 & 10.6832\\
APR4 & gradient & 10000 & 2.248766 & 1.370089 & 76.277 & 5.327196 & 2.141379 & 10.7297\\
APR4 & gradient & 100000 & 2.198771 & 1.379310 & 75.753 & 5.243779 & 2.090335 & 10.6985\\
\midrule
SLy & mass & 1000 & 2.049739 & 1.000559 & 99.934 & 0.237140 & 2.049739 & 9.9802\\
SLy & mass & 10000 & 2.098219 & 1.011216 & 98.689 & 1.067392 & 2.098073 & 10.2392\\
SLy & mass & 100000 & 2.053290 & 1.000523 & 99.939 & 0.229765 & 2.053289 & 10.0132\\
\bottomrule
\end{tabular}
\end{table}

\subsection{Turning-point heuristic}
On the 53-point scan grid every sequence has exactly one sign change of $dM/d\rho_c$, at its mass maximum. The
branch with $dM/d\rho_c>0$ runs from the grid floor $\log_{10}\rho_{0c}=14.55$ ($M=0.23\,\Msun$ APR4,
$0.29\,\Msun$ SLy, $0.70\,\Msun$ H4; lower densities were not computed) to the refined maximum, and the 1.4 and
$2.0\,\Msun$ points are on that branch in all 33 sequences. When the optimizer and root-finder points are
included, additional sign flips appear; all lie within $\pm\ee{2}{-4}$~dex of the peak, with reversal steps
$\le\ee{2.1}{-7}\,\Msun$, below the $\ee{6.4}{-7}\,\Msun$ refinement floor (Sect.~\ref{sec:conv}). They are
numerical noise on a flat maximum, not additional turning points; the peak location itself is resolved only
to about $10^{-4}$~dex (\file{stellar_summary.csv}; \file{causality_stability_log.txt}). The lower end of this
branch is the grid limit, not a physical stability boundary. For the modified equations the result is a
heuristic, not established radial stability.

\subsection{Illustrative NICER overlay}
Figure~\ref{fig:mr} shows the 33 sequences with the three plotted NICER references of Table~\ref{tab:nicer}.
TAL curves are solid and GR curves dashed up to the causal endpoint; the continuation to the sequence maximum is
dotted and faded and is a potentially acausal benchmark of the published fit. Open circles mark TAL and open
squares GR causal endpoints; markers appear only on sequences that cross the onset (APR4: 11 of 11; SLy: 3 of
11; H4: none). The overlay is a visual comparison with approximate ellipses of GR-derived posteriors; it is not a
TAL pulse-profile likelihood or a Bayes factor, and rotation, which changes equatorial radius and maximum mass,
is not modelled.

\begin{sidewaysfigure}
\centering
\includegraphics[width=\textheight]{figures/mr_nicer.pdf}
\caption[]{Conditional mass--radius sequences for SLy, APR4 and H4 (Read fits), top row mass-source ansatz,
bottom row gradient ansatz, one column per $L$. Solid: TAL; dashed: GR, each up to the causal endpoint;
dotted, faded: potentially acausal continuation to the sequence maximum. Shading: independent, symmetrized
Gaussian joint 68\,\%/95\,\% approximations of the GR-derived NICER results (not published posterior
contours, not a TAL likelihood). Horizontal line: nominal $2\,\Msun$ mass-capability check. Modified-branch
radial stability is not established. Package file \file{mr_nicer.pdf}, reproduced unchanged.}
\label{fig:mr}
\end{sidewaysfigure}

\subsection{Sirius B}
With $R_\odot=\ee{6.957}{8}$~m, $R_B=\ee{5.586471}{6}$~m and the compactness is $\ee{2.69079}{-4}$. The exact
exterior Schwarzschild expression gives $v_{\mathrm{GR}}=80.70055$~km\,s$^{-1}$ with $\sigma_{\mathrm{GR}}=1.4080$~km\,s$^{-1}$;
the measured-minus-GR residual is $-0.05055\pm1.60481$~km\,s$^{-1}$ after adding the redshift measurement
error. The approximate two-sided-95\,\% positive extra-shift allowance is 3.0949~km\,s$^{-1}$; with the
larger quoted redshift-mass uncertainty it is 4.1645~km\,s$^{-1}$ (\file{wd_summary.json}). Both are approximate
Gaussian allowances, not a covariance-aware posterior on $L$; mass and radius share astrometric inputs.

At $L=2R_B=\ee{1.1172942}{7}$~m the all-enclosed formula adds approximately 9.63~km\,s$^{-1}$, larger than either
allowance. Holding the GR exterior and the dynamical interpretation fixed, the profile-independent conditional
exclusion is $\ee{1.117}{7}\lesssim L\lesssim\ee{3.477}{7}$~m with direct independent-error propagation, and
$\ee{1.117}{7}\lesssim L\lesssim\ee{2.584}{7}$~m with the more conservative quoted-mass allowance
(\file{numerical_diagnostics.json}, \file{profile_independent_WD_band}). These are a subset of the possible
excluded band under this specific additive-clock model, not global bounds on $L$.

\begin{figure}[t]
\centering
\includegraphics[width=0.85\textwidth]{figures/sirius_redshift.png}
\caption[]{Sirius~B: additional differential redshift under the conditional additive-clock coupling with the GR
exterior retained, for WD polytropes $n=1.5$ and $n=3$ (Sirius~A $n=3$), versus $L$. Dashed: approximate
95\,\% extra-shift upper allowance. The curve falls as $1/L$ once the whole white dwarf is enclosed, so no
monotonic upper limit on $L$ follows. Package file \file{sirius_redshift.png} (data:
\file{sirius_redshift.csv}), reproduced unchanged.}
\label{fig:sirius}
\end{figure}

For the illustrative profiles (Table~\ref{tab:wd}), the sampled peak extra shifts are about 11.56~km\,s$^{-1}$
at $L\approx\ee{8.25}{6}$~m for $n=1.5$ and 13.57~km\,s$^{-1}$ at $L\approx\ee{7.26}{6}$~m for $n=3$
(Fig.~\ref{fig:sirius}); these are dimensionless clock corrections of order $\ee{4}{-5}$, not percent-level
factors. At 100~km even the larger illustrative shift is only 0.0076~m\,s$^{-1}$, compared with the measured
770~m\,s$^{-1}$ redshift error. The sampled curves therefore do not constrain any of the $L\le100$~km values,
and once the whole white dwarf is enclosed the extra shift decreases as $1/L$: there is no monotonic upper
limit on $L$ from this finite-source argument. The very small-$L$ entries are idealized zero-density-surface
numbers, not atmosphere-resolved predictions.

\begin{table}[t]
\centering
\caption[]{Extra surface redshift of Sirius~B for the illustrative polytropes. Sources: \file{sirius_redshift.csv};
\file{note_HU_EN.md} \S5.}
\label{tab:wd}
\begin{tabular}{rcc}
\toprule
$L$ (m) & $\delta v_B$ (km\,s$^{-1}$), $n=1.5$ & $\delta v_B$ (km\,s$^{-1}$), $n=3$\\
\midrule
1 & $\ee{2.4}{-23}$ & $\ee{3.1}{-34}$\\
367 & $\ee{2.3}{-14}$ & $\ee{2.0}{-21}$\\
1000 & $\ee{7.6}{-13}$ & $\ee{3.0}{-19}$\\
10000 & $\ee{2.4}{-9}$ & $\ee{3.0}{-14}$\\
100000 & $\ee{7.6}{-6}$ & $\ee{3.1}{-9}$\\
\bottomrule
\end{tabular}
\end{table}

\subsection{Convergence}\label{sec:conv}
Over the 66 refinement comparisons the largest differences were $\ee{6.43}{-7}\,\Msun$ in mass and
$\ee{3.92}{-6}$~km in radius; the estimated fixed-mass radius difference near $1.4\,\Msun$ was at most
$\ee{4.10}{-6}$~km, and the refined exponent fixed-point residual at most $\ee{4.41}{-8}$
(\file{convergence.csv}). For the causal endpoints and core boundaries the maximum absolute differences were
$\ee{8.5}{-7}\,\Msun$ (endpoint mass), $\ee{3.8}{-6}$~km (endpoint radius) and $\ee{7.0}{-7}$~km (core
boundary) (\file{causal_endpoint_convergence.csv}). The exact-integral cross-check agrees with the production
convolution to within $\ee{9.0}{-5}$ relative in all 15 cases (\file{numerical_diagnostics.json}). These are
tested changes, not rigorous global error bounds; numerical precision is not EOS accuracy, and the Read fits
themselves differ from the microscopic tabulations.

\subsection{Reproduction of the V21 run}
The scripts were run as five commands: three independent \file{--eos} runs, then \file{--combine}, then
\file{diagnostics.py} with Python 3.13.5, NumPy 2.5.4, SciPy 1.18.1, pandas 3.1.0 and matplotlib 3.11.2
(\file{numerical_diagnostics.json}; \file{README.txt}). Relative to the earlier run (Python 3.12.3, pandas 3.0.6)
the outputs are numerically, not byte-for-byte, reproduced: relative $M_{\max}$ change $\le\ee{1.7}{-7}$,
fixed-mass radii $\le\ee{1.2}{-7}$ relative ($\le\ee{1.4}{-6}$~km), radius shifts $\le\ee{1.8}{-6}$~km. The peak
central density moved by $\le\ee{1.4}{-4}$~dex on the flat maximum, so the radius at $M_{\max}$ moved by
$\le\ee{9.4}{-4}$~km, and \file{mr_sequences.csv} has 3113 rather than 3103 rows. All headline numbers are
unchanged at their quoted precision (\file{note_HU_EN.md} \S7.3).

\section{Discussion}
\paragraph{Status of the results.} Within the stated conditional model, a locally centred top-hat produces
percent-level clock factors in neutron stars but not in white dwarfs or the Sun. The resulting radius shifts at
fixed mass are smaller than the spread of three standard EOS fits at 1.4 and $2.0\,\Msun$. These are sensitivity
calculations, not detections, and they do not establish a closed physical consistency of TAL.

\paragraph{Limitations and open items.}
\begin{enumerate}[leftmargin=*,itemsep=2pt]
\item \emph{Coupling consistency (task 50aa2db0).} Open. The two ans\"atze are not derived from a complete TAL
field equation; task 50aa2db0 may change which insertion, if any, is admissible.
\item \emph{Radial stability.} The $dM/d\rho_c>0$ test is a heuristic. The perturbation operator, boundary
conditions and mode frequencies of the modified equations have not been computed; the \file{stable_branch_*}
fields must be read with this restriction.
\item \emph{Causal EOS completion.} Not done. APR4 is superluminal above $\rho_0=\ee{1.4458}{15}$~g\,cm$^{-3}$ in
GR and in all modified branches; SLy marginally in three mass-source maxima. \file{M_max_causal_Msun} is the
endpoint of the fit's causal range, not the maximum mass of a causally completed EOS.
\item \emph{NICER inference.} The overlay uses approximate ellipses of GR-derived posteriors; no TAL
pulse-profile likelihood, Bayes factor or rotating stellar model was computed.
\item \emph{Low densities.} The branch below $\log_{10}\rho_{0c}=14.55$ is not scanned.
\item \emph{Sirius~B.} The additive clock coupling and the retained GR exterior are separate assumptions; the V21
causality extension does not establish them. Correlated astrometric errors, atmospheric systematics or a
replacement rather than additive law could change the inferred band.
\item \emph{Model choices.} Which density enters TAL ($\epsilon/c^2$ versus $\rho_0$), Euclidean versus
proper-volume averaging, and the mass-source versus gradient coupling are modelling choices, not uncertainty
bars.
\end{enumerate}

\section{Conclusion}
For three piecewise-polytropic EOS fits and five averaging radii, the TAL factor $f_L=\exp(-\gamma_L\bar\rho_L)$
with a locally centred, self-consistently recomputed top-hat was inserted into the stellar equilibrium equations
in two labelled ways. The largest fixed-mass radius shift is 0.164~km at $1.4\,\Msun$ (three-EOS GR spread
2.632~km) and 0.623~km at $2.0\,\Msun$ (spread 1.557~km), all in the causal range of the fits; radius alone
therefore does not identify TAL in this comparison. APR4 maxima lie above its acausal onset already in GR and are
reported only as labelled benchmarks, with causal endpoints given separately. Under a separately assumed
additive clock coupling, Sirius~B excludes only a finite stellar-scale band of $L$ and leaves the 1~m--100~km
values unconstrained. Task 50aa2db0, a full radial-mode stability analysis, a causal EOS completion and a TAL
NICER inference remain open. No detection claim or unconditional bound on $L$ is made.

\section*{Data availability}
All inputs, outputs and the note are archived in Zenodo record V21, ``Time Attraction Law V21 --- final
academy-verified package'', doi:10.5281/zenodo.23602501, published 2026-10-23 \cite{v21} (files
\file{ns_tal_V21_2026-10-23 (3).zip} and \file{note_HU_EN (3).md}). The identical package is GitHub release V19,
directory \file{V21/} of \url{https://github.com/unbekanntedavid-glitch/Time-Attraction}, commit
\texttt{\seqsplit{0478d663a0e4a89afdb69dc57b6f76bdc0c8af94}} \cite{gh}; the package zip \file{ns_tal_V21_2026-10-23.zip} has SHA-256
\texttt{\seqsplit{dc1768aa8148fe989d142256ed8c27518042ed6fecf09b4cbfd0dd68e43de56d}}, and \file{SHA256SUMS.txt} lists the 19
package files. The previous releases are Zenodo V20 (GitHub V18), doi:10.5281/zenodo.23540741 \cite{v20}, and V19,
doi:10.5281/zenodo.23497128 \cite{v19}. All materials are licensed under CC BY-NC 4.0.

\section*{Code availability}
The calculation is reproduced by \file{stellar_tal.py} and \file{diagnostics.py} in the package, written in
Python~3 with NumPy \cite{numpy}, SciPy \cite{scipy}, pandas and matplotlib. No network access, external EOS library or fitted
TAL parameter is needed. The commands are
\file{python3 stellar_tal.py --eos SLy}, \file{--eos APR4}, \file{--eos H4}, then
\file{python3 stellar_tal.py --combine} and \file{python3 diagnostics.py} (\file{README.txt}). Figures~\ref{fig:mr}
and~\ref{fig:sirius} are the package figures, included unchanged.

\section*{Declarations}
\textbf{Funding.} No funding was received for this work.
\textbf{Competing interests.} The author has no competing interests to declare.
\textbf{Author contributions.} D.\,Sz.\ is the sole author.
\textbf{Ethics approval, consent.} Not applicable.
\textbf{Licence.} CC BY-NC 4.0; all rights, including commercial rights, remain with the author.

\appendix
\section{Audit fixes and archive bookkeeping}
The V21 package was audited before release (\file{note_HU_EN.md} \S7.4). The fixes were:
H1, figure styles and causal endpoint markers as described for Fig.~\ref{fig:mr};
H2, the core radius taken as the areal radius of the $c_s^2=c^2$ boundary from the dense ODE solution, with the
previous sampled value retained in separate columns as a lower estimate (it is exceeded by up to 10.68~m at
$M_{\max}$ and 17.6~m over all 654 acausal sequence points);
H3, \file{plot_label} used in the legend and caption and the fail-loud symmetrized-error check;
H4, five executable commands in \file{README.txt}; and the endpoint/core-boundary refinement tests of
Sect.~\ref{sec:conv}.

Archive count of acausal points: selecting the \file{acausal_point=True} rows of \file{mr_sequences.csv} whose
\file{log10_rho0c_cgs} is less than or equal to the peak central density stored in \file{stellar_summary.csv}
for the same EOS, coupling and $L$, compared as exact decimal values without additional rounding or tolerance,
gives 182 unique rows (173 APR4, 9 SLy). This is an archive bookkeeping count; repeated appearances across figure
panels are not counted again. The clarification of this rule changed only \file{note_HU_EN.md},
\file{README.txt} and \file{SHA256SUMS.txt}; code, numerical tables, figures and conclusions are unchanged.

\begin{thebibliography}{99}\small
\bibitem{will} C. M. Will, The confrontation between general relativity and experiment, Living Rev. Relativ. 17, 4
(2014). doi:10.12942/lrr-2014-4
\bibitem{tal} D. Szabolcsi, Az Id\H{o}vonz\'as T\"orv\'enye: $f_L=\exp(-\gamma_L\bar\rho_L)$ [The law of time
attraction: density-dependent local time scale], Zenodo (2026), file \file{Time_Attraction.pdf}.
doi:10.5281/zenodo.23066993
\bibitem{v20} D. Szabolcsi, Time Attraction Law V20 --- audited release with independent third-round review,
Zenodo, version v20 (2026). doi:10.5281/zenodo.23540741
\bibitem{v19} D. Szabolcsi, Time-Attraction Law seventh model independent verification, Zenodo (2026).
doi:10.5281/zenodo.23497128
\bibitem{v21} D. Szabolcsi, Time Attraction Law V21 --- final academy-verified package, Zenodo (2026).
doi:10.5281/zenodo.23602501
\bibitem{gh} D. Szabolcsi, Time-Attraction, GitHub repository, commit \texttt{\seqsplit{0478d663a0e4a89afdb69dc57b6f76bdc0c8af94}}
(directories V21/ and V18/, 2026). \url{https://github.com/unbekanntedavid-glitch/Time-Attraction}
\bibitem{read} J. S. Read, B. D. Lackey, B. J. Owen, J. L. Friedman, Constraints on a phenomenologically
parametrized neutron-star equation of state, Phys. Rev. D 79, 124032 (2009). doi:10.1103/PhysRevD.79.124032
\bibitem{salmi} T. Salmi et al., The Radius of the High-mass Pulsar PSR J0740+6620 with 3.6 yr of NICER Data,
Astrophys. J. 974, 294 (2024). doi:10.3847/1538-4357/ad5f1f
\bibitem{choudhury} D. Choudhury et al., A NICER View of the Nearest and Brightest Millisecond Pulsar:
PSR J0437$-$4715, Astrophys. J. Lett. 971, L20 (2024). doi:10.3847/2041-8213/ad5a6f
\bibitem{kini} Y. Kini et al., A NICER View of PSR J0030+0451: Updated Constraints from Six Years of NICER
Observations (2026). arXiv:2602.23743
\bibitem{miller} M. C. Miller et al., The Radius of PSR J0437$-$4715 from NICER Data (2025). arXiv:2512.08790
\bibitem{joyce} S. R. G. Joyce et al., The gravitational redshift of Sirius B, Mon. Not. R. Astron. Soc. 481,
2361--2370 (2018). doi:10.1093/mnras/sty2404
\bibitem{atnf} R. N. Manchester, G. B. Hobbs, A. Teoh, M. Hobbs, The ATNF Pulsar Catalogue, Astron. J. 129, 1993
(2005); catalogue v2.6.5, \url{https://www.atnf.csiro.au/research/pulsar/psrcat/}
\bibitem{bond} H. E. Bond et al., The Sirius System and Its Astrophysical Puzzles: Hubble Space Telescope and
Ground-based Astrometry (2017). doi:10.3847/1538-4357/aa6af8; arXiv:1703.10625
\bibitem{numpy} C. R. Harris et al., Array programming with NumPy, Nature 585, 357--362 (2020).
doi:10.1038/s41586-020-2649-2
\bibitem{scipy} P. Virtanen et al., SciPy 1.0: fundamental algorithms for scientific computing in Python, Nat.
Methods 17, 261--272 (2020). doi:10.1038/s41592-019-0686-2
\end{thebibliography}
\end{document}
